\documentclass[10pt,a4paper]{amsart}
\usepackage[margin=19mm]{geometry}
\usepackage{amsmath,amssymb,mathtools}
\usepackage[scr=rsfs,cal=euler]{mathalfa}
\usepackage{xfrac}
\usepackage[unicode,hidelinks,bookmarks=false]{hyperref}
\newcommand{\ie}{i.e.\ }
\newcommand{\cf}{cf.\ }
\newcommand{\resp}{respectively\ }
\newcommand{\Subsection}{\subsection}
\providecommand{\nobreakdash}{\nobreak\hbox{-}}
\setlength{\emergencystretch}{2em}
\multlinegap0pt
\newtheorem{thm}{Theorem}
\newtheorem{lem}{Lemma}
\def\fish{\mbox{\rm Fish}}
\title{Constraint satisfaction problems, compactness and non-measurable sets}
\author{Claude Tardif}
\date{Source: arXiv:2508.14838v4 (CC BY 4.0)}
\begin{document}
\maketitle

\noindent\emph{Source and licence.} Constraint satisfaction problems, compactness and non-measurable sets. Version arXiv:2508.14838v4 (CC BY 4.0), distributed by arXiv under the Creative Commons Attribution 4.0 International licence, http://creativecommons.org/licenses/by/4.0/, as recorded in the arXiv licence metadata for this exact version and shown on the abstract page https://arxiv.org/abs/2508.14838v4. Full licence notice: \texttt{licenses/arxiv-2508.14838-CC-BY-4.0.txt}.\par\smallskip

\noindent\textit{Editorial scope.} Counted unit: the article's thm main (source0.tex lines 100--104). The article's complete original proof of it is delivered with it (source0.tex lines 369--389, one \texttt{proof} environment of 21 lines, which the article places away from the statement). No mathematics is written, repaired or re-derived: the statement and the proof are the article's own lines, in the article's own notation, and no retained line is substituted or re-flowed.  Retained uncounted as the source-local context the proof needs: lines 317--320; lines 333--335.  Nothing was omitted from any retained span, so the package carries no omission record.  No retained line contains a citation, so the wrapper carries no bibliography and no bibliography entry.  No retained line contains a figure environment, an image-inclusion command or drawing code, so no image is delivered and the package declares no asset dependency.  Editorial formatting only, in addition: an AMS \texttt{amsart} preamble that restates the article's own declarations and macros used by the retained lines, namely the environments \texttt{thm}, \texttt{lem} on separate counters, together with the standard packages those lines need.  The retained environments print in this document as Theorem 1, Lemma 1 in order of appearance, whereas the article prints the same items with the numbering of its own full document.  The complete native source of the article as distributed in the arXiv e-print for this exact version is bundled unchanged as \texttt{sources/183577/source0.tex}.
\begin{lem} \label{compacthom}
If $A$ is compact, then there exists a homomorphism $\phi: C \rightarrow A$.
\hfill $\qed$
\end{lem}

\begin{lem}\label{mezero}
If $\fish_{f,D}$ is measurable, then $\mu(\fish_{f,D}) = 0$.
\end{lem}

\begin{thm}\label{main}
Let $A$ be a relational structure. If $A$ does not have width $1$, then the axiom ``$A$ is compact''
implies the existence of a set in $\mathbb{R}^3$
that is not Lebesgue measurable.
\end{thm}

\begin{proof}[Proof of Theorem~\ref{main}] 
Suppose that $A$ is compact and all sets in $\mathbb{R}^3$
are measurable. Then by Lemma~\ref{compacthom},
there exists a homomorphism $\phi: C \rightarrow A$.
This allows to define the functions $f_p$, 
$p \in \mathbb{S}_2\setminus F_\mathcal{G}$
and the sets $\fish_{f,D}$, $f \in V(A)^{V(\mathcal{U}(A))}$, 
$D \in \mathcal{D}$.
By Lemma~\ref{mezero}, these sets all have measure $0$,
and so does their finite union.
The set $\{\lambda p | \lambda \in (0,1], p \in F_{\mathcal{G}}\}$
also has measure $0$ since it can constructively be contained
in a countable collection of boxes of arbitrarily small measure.
Therefore there exists a point $p$ of $\mathbb{S}_2$ 
such that the line $\{\lambda p | \lambda \in (0,1]\}$ 
belongs to neither of these sets.
We then have $f_p = f_{D p}$ for all $D \in \mathcal{D}$. This implies that $f_p$ is a homomorphism
from $\mathcal{U}(A)$ to $A$, since collapsing each set 
$\{ (D p/\mathcal{N},S) | D \in \mathcal{D} \}$ to a single 
element yields a structure isomorphic to $\mathcal{U}(A)$.
\end{proof}

\end{document}
